\subsection{Positive Oracle Estimates from Identical Programs}
\label{sec:f3-clones}
All nine clone slots implement the same program: their population-level
stable complementarity is zero. On the common $\RcvPairedN$ tasks,
their repeat-averaged oracle estimate is $\RcvClonePlugin$ pp, while
the generated panel gives $\RcvRealPlugin$ pp
(Table~\ref{tab:clone-control}). Both estimates use identical task
identities and complete observations. The control establishes that
averaging three repeats does not make a positive oracle estimate
identify specialization.

\begin{table}[htbp]
\centering\small
\begin{tabular}{@{}lcc@{}}
\toprule
Quantity & Generated + \bare{} & Same-code clones \\
\midrule
Complete-task oracle estimate $\widehat H$ & $\RcvRealPlugin$ pp & $\RcvClonePlugin$ pp \\
Complete tasks for repeat-based $G$ & \RcvRealN & \RcvCloneN \\
Held-out-repeat mean gain $G$ & $\RcvRealG$ pp & $\RcvCloneG$ pp \\
Residual score-pattern correlation & $\RcvRealRho$ & $\RcvCloneRho$ \\
\midrule
\multicolumn{3}{l}{Paired clone-adjusted mean: $D=\RcvD$ pp on $\RcvPairedN$ tasks} \\
\multicolumn{3}{l}{Analysis-set completeness: $100\%$ in both arms ($\RcvCells$ cells each)} \\
\bottomrule
\end{tabular}
\caption{\textbf{Coverage, repeatability, and selection gain separate.}
All quantities use the same $\RcvPairedN$ tasks. Both arms are fully
observed on this analysis set. $D$ is the paired difference of the
unrounded cross-fitted gains; rounding may affect the displayed difference.
}
\label{tab:clone-control}
\end{table}

\subsection{Repeatable Score Patterns and Held-Out Selection Gain}
\label{sec:repeat-result}
Generated programs do exhibit repeatable differences. An exploratory
analysis removes additive member and task effects from the score-advantage
matrix and correlates each held-out repeat with the other two.
The mean residual correlation is $\RcvRealRho$ for generated programs
and $\RcvCloneRho$ for clones; their paired-bootstrap difference is
$\RcvRhoDiff$ (95\% interval $\RcvRhoDiffCI$).
These correlations describe repeatable \emph{scored-outcome patterns}.
To examine their direction, we count persistent losses and wins.
All-three-repeat losses to
\bare{} span $\RcvLossTasks$ tasks ($\RcvLossPairs$ member--task pairs), while all-three-repeat
wins occur on only one task, shared by all eight generated members.
That apparent repair is sensitive to answer-format extraction
(Appendix~\ref{app:repeatability}). The observed asymmetry helps explain
the coexistence of strong score repeatability and limited observed wins.

The held-out selection diagnostic asks the stronger question: do past
outcomes identify a useful member for another execution?
The clone-adjusted mean is $D=\RcvD$ pp, with 95\% interval
$\RcvDCI$ pp and randomization $p=\RcvP$.
All primary analysis tasks are fully observed. The decision nevertheless
remains \textsc{Insufficient Evidence}: the interval includes zero and
the positive-support conditions are not met; the result does not
establish equivalence.
Excluding the three included tasks containing conflicting scores gives
$D=\RcvConflictD$ pp
on $\RcvConflictN$ paired tasks and the same decision
(Appendix~\ref{app:wp1r}).

\paragraph{Would the diagnostic detect real crossovers?}
\label{sec:power}
Known-truth calibration quantifies the diagnostic's sensitivity. Across the tested $100$-cell grid the worst observed
null false-support rate is $2\%$, but a two-specialist setting with
true $H_{\rm stable}=17.5$ pp has $0\%$ power at $R=3$, about $0.7\%$
at $R=5$, and $100\%$ at $R=10,20$. Smaller cyclic crossovers remain
largely undetected even at $R=20$ (Appendix Table~\ref{tab:calibration}).
These are setting-specific empirical results. Thus the data identify
repeatable score patterns, while useful task-wise crossovers remain
unresolved. Actual pre-execution selection can still be evaluated directly.
